\begin{frame}{Conclusion}

    \begin{itemize}
        \item A well-defined information standard facilitates both organization-wide and inter-organization data management
        \begin{itemize}
            \item It allows the use of expert models for \textit{every} domain
            \item Standardization of information helps with every aspect of (FAIR) data management
        \end{itemize}
        \item Defining and standardizing information is about standardizing standards
        \begin{itemize}
            \item The motivations are identical
            \item The methods are identical
            \item Just one level of abstraction higher
        \end{itemize}
        \item A well-defined information standard requires
        \begin{itemize}

            \item A clear distinction between vocabulary models and constraint models (req. 1)
            \item Harmonized metamodels; choosing an appropriate pivot model is key (req. 2)
            \item A fair amount of interpretation, but an upper ontology can provide a structure for interpretation (req. 3)
        \end{itemize}
        \item Applied to openADR 3.0: what the standard makes normative is the \emph{serialization} of the model
        — the OpenAPI document — and the requirements of the day apply to that artifact directly
        %\item Most of what I have discussed today is at best in the experimental phase, so there is a lot of work to be done \pause (and yes, I think EPRI might have a role in that)
    \end{itemize}
\end{frame}
